\subsection{Pin 结构}\label{sec:11i}
$\operatorname{Pin}_n$ 是 $O_n$ 关于 $\Z/2$ 的中心扩张。这样的扩张实际有两个，它们拓扑相同，却是不同的 Lie 群。对我们的目的，其中一种稍方便，因此规定 $\operatorname{Pin}_n$ 为由 $w_2\in H^2(BO_n;\Z/2)$ 在拓扑群扩张意义下分类的群。\znote{167}{两种约定源于对二次型取 $q_\pm=\pm\sum_i(x^i)^2$：$\operatorname{Pin}_+(n)$ 满足 $v^2=+\|v\|^2$，$\operatorname{Pin}_-(n)$ 满足 $v^2=-\|v\|^2$，两者都是 $O_n$ 的二重覆盖，因而是同胚的拓扑空间，却有不同的群结构。正文取 $v^2=\|v\|^2e$，故这里的 $\operatorname{Pin}_n$ 即 $\operatorname{Pin}_{0,n}$，也记作 $\operatorname{Pin}_n^+$。这一约定还决定了反射的提升是二阶还是四阶元素：$\operatorname{Pin}_+$ 中反射的原像满足 $a^2=e$，而另一版本 $\operatorname{Pin}_{n,0}=\operatorname{Pin}_n^-$ 中它满足 $a^4=e$，且其分类类为 $w_2+w_1^2$，正如下一则笔记所述。两个版本给出的 Pin 结构障碍因此相差 $w_1^2$，这正是选用哪一版本会改变障碍计算的原因。[来源：nLab "pin group"，https:/\allowbreak /\allowbreak ncatlab.org/\allowbreak nlab/\allowbreak show/\allowbreak Pin+\allowbreak group（2026-11-15 访问）]}简要回顾其具体定义。将 $\R^n$ 嵌入 Clifford 代数 $\operatorname{Cl}(\R^n)$，其中向量满足 $v^2=\|v\|^2e$，$e\in\operatorname{Cl}(\R^n)$ 为单位元。$\operatorname{Pin}_n\subset\operatorname{Cl}(\R^n)^\times$ 是由 $S^{n-1}\subset\R^n$ 生成的乘法子群。\znote{167}{nLab 的 pin 群条目给出的定义与此一致：$\operatorname{Pin}(V,q)$ 是 Clifford 代数乘法群 $\operatorname{GL}_1(\operatorname{Cl}(V,q))$ 的子群，由满足 $q(v_i)=1$ 的向量之积 $v_1\cdots v_n$ 生成；其中偶数个因子之积再构成 $\operatorname{Spin}(V,q)$。作为抽象群，它由这些单位向量生成，关系来自 Clifford 关系 $v^2=q(v)$；投影 $\operatorname{Pin}(V,q)\to O(V)$ 由扭曲伴随作用给出，核恰为 $\{\pm1\}$。正文中的 $v$ 满足 $v^2=\|v\|^2e$，即 $q(v)=1$，故正是这里的生成元。这些关系在 $n$ 固定后是有限的，故 $\operatorname{Pin}(V,q)$ 是有限维 Lie 群，维数与 $O(V)$ 相同。标准参考是 Lawson–Michelsohn《Spin geometry》第 I 章第 2 节，nLab 该条目也列出了此书。[来源：nLab "pin group"，https:/\allowbreak /\allowbreak ncatlab.org/\allowbreak nlab/\allowbreak show/\allowbreak Pin+\allowbreak group（2026-11-15 访问）]}同态 $\operatorname{Pin}_n\to O_n$ 是扭曲伴随作用在 $\R^n\subset\operatorname{Cl}(\R^n)$ 上的限制：元素 $v\in S^{n-1}\subset\operatorname{Pin}_n$ 按 $w\mapsto-vwv$ 作用，即关于垂直于 $v$ 的超平面反射。当然，可以用任意内积空间 $F$ 代替 $\R^n$，定义具有同样性质的群 $\operatorname{Pin}(F)$。考察反射的复合可得：
\begin{lemma}\label{st:11.14}
若 $A\in O(F)$ 为对合，则 $A$ 在 $\operatorname{Pin}(F)$ 中的原像与子空间 $F^{\mathrm{anti}}=\ker(\Id+A)$ 的定向典范对应。任一这样的原像的平方为
\begin{equation}\label{eq:11.29}
(-1)^{\frac12\operatorname{rank}(F^{\mathrm{anti}})(\operatorname{rank}(F^{\mathrm{anti}})-1)}e.
\end{equation}
特别地，反射的原像是二阶元素。
\end{lemma}
\begin{note}\label{st:11.15}
Pin 的一般参考文献有\bcite{93}，其中我们的 $\operatorname{Pin}_n$ 被称为 $\operatorname{Pin}_{0,\allowbreak n}$；以及\bcite{86}，其中称它为 $\operatorname{Pin}_n^+$。另一个版本记为 $\operatorname{Pin}_{n,\allowbreak 0}$ 或 $\operatorname{Pin}_n^-$，具有以下性质：它关联于类 $w_2+w_1^2$；它位于
\SourcePage{175}
满足 $v^2=-\|v\|^2e$ 的、更为熟知的 Clifford 代数内；最后，在该群中反射的原像是四阶元素。
\end{note}
\begin{definition*}
如通常一样，实内积向量丛 $F^n\to B$ 上的 \emph{Pin 结构}，是一个 $\operatorname{Pin}_n$ 主丛 $P^\#\to B$，连同同构 $P^\#\times_{\operatorname{Pin}_n}\R^n\cong F$。
\end{definition*}
同态 $\operatorname{Pin}_n\to O_n$ 诱导二重覆盖 $P^\#\to P$，其中 $P$ 为 $F$ 的标准正交标架丛。由最初对 $\operatorname{Pin}_n$ 的刻画，这种结构存在的障碍是 $w_2(F)$。若该类消失，则 Pin 结构的同构类所成的集合是 $H^1(B;\Z/2)$ 上的仿射空间。\znote{168}{Pin 结构的标准图景是正交标架丛分类映射的提升：标架丛给出映射 $B\to BO_n$，Pin 结构即它穿过 $\operatorname{BPin}_n\to BO_n$ 的提升，二重覆盖的障碍正是 $w_2$。若两个提升存在，它们的差别给出一个到 $K(\Z/2,1)$ 的映射，即 $H^1(B;\Z/2)$ 中的一个类；反过来，任给这样的类都可用它扭出新的提升，所以全部 Pin 结构构成 $H^1(B;\Z/2)$ 的仿射空间。nLab 关于扭曲 spin 结构的条目以同样方式刻画了 spin 结构为 Cech 圈穿过 $\mathrm{BSpin}\to BSO$ 的提升。[来源：nLab "twisted spin structure"，https:/\allowbreak /\allowbreak ncatlab.org/\allowbreak nlab/\allowbreak show/\allowbreak twisted+\allowbreak spin+\allowbreak structure（2026-11-15 访问）]}具体地，给定 Pin 结构 $P^\#$ 和实线丛 $\beta\to B$，可以取纤维积 $P^\#\times_{\Z/2}S(\beta)$ 形成新的 Pin 结构；这里 $\Z/2$ 以 $\pm1$ 作用于与 $\beta$ 关联的二重覆盖 $S(\beta)$，并通过中心子群 $\ker(\operatorname{Pin}_n\to O_n)=\{\pm e\}$ 作用于 $P^\#$。为简便，将新结构记为 $P^\#\otimes\beta$。

任意向量丛 $F$ 都可关联一个群丛 $\operatorname{Pin}(F)$。下面澄清一个可能混淆的地方。给定 $F$ 上的 Pin 结构 $P^\#$，考虑群丛 $\operatorname{Aut}(P^\#)\to B$，其在 $x$ 上的纤维由 $\operatorname{Pin}_n$ 主齐性空间 $P_x^\#$ 的全部自同构组成。可将其视为伴随丛
\[
\operatorname{Aut}(P^\#)=P^\#\times_{\operatorname{Pin}_n}\operatorname{Pin}_n
=P\times_{O_n}\operatorname{Pin}_n\hookrightarrow P\times_{O_n}\operatorname{Cl}(\R^n).
\]
这里 $\operatorname{Pin}_n$ 通过共轭作用于自身；此作用下降为 $O_n$ 作用，并嵌入 $\operatorname{Cl}(\R^n)$ 上的非扭曲伴随作用。例如，关于超平面 $v^\perp$ 的反射按 $w\mapsto vwv$ 作用，其在 $\R^n\subset\operatorname{Cl}(\R^n)$ 上的限制等于该反射乘以 $-1$。所以 $\operatorname{Aut}(P^\#)$ 同构的并非 $\operatorname{Pin}(F)$，而是 $\operatorname{Pin}(F\otimes\lambda^{\mathrm{top}}(F))$，并构成如下交换图：
\SourceEquationDiagram{diagram-11-30}{0.51507637}{11.30}
更一般地，将 Pin 结构 $P^\#$ 分别与两个线丛 $\beta_0,\allowbreak \beta_1$ 作张量积，所得主丛之间的 $\operatorname{Pin}_n$ 等变映射所成的丛为
\begin{equation}\label{eq:11.31}
\operatorname{Iso}(P^\#\otimes\beta_0,P^\#\otimes\beta_1)
\cong\operatorname{Pin}(F\otimes\lambda^{\mathrm{top}}(F))\times_{\Z/2}S(\beta_0^\vee\otimes\beta_1).
\end{equation}
\begin{note}\label{st:11.16}
还有关于参考类 $w\in H^2(B;\Z/2)$ 扭曲的 Pin 结构的一般理论，其存在障碍为 $w_2(F)+w$。一个 gerbe 式定义如下：选择分别表示 $w$ 与丛 $F$ 的 Čech 上链 $\check w\in\check C^2(B;\Z/2)$ 与 $\check F\in\check C^1(B;O_n)$。扭曲 Pin 结构
\SourcePage{176}
是 $\check C^1(B;\operatorname{Pin}_n)$ 中投影到 $\check F$ 的元素，而它未能成为上闭链的偏差由 $\check w$ 编码；因为 $\Z/2\subset\operatorname{Pin}_n$ 是中心子群，所以这一定义有意义。\znote{169}{扭曲 spin 结构的一般定义是：预先指定障碍类 $\alpha\in H^2(B;\Z/2)$，要求标架丛的 Čech 圈 $g:B\to BO_n$ 满足 $[w_2(g)]=\alpha$；所有这样的圈连同其间的同伦构成一个 groupoid，它由 $(\infty,1)$-拉回刻画。障碍本身是一个 $\mathrm{B}\Z/2$ 主 $2$-丛，即束 gerbe（Murray–Singer 称之为 spin gerbe）。正文这里的 gerbe 式定义正是该说法的 Čech 版本，其中 $\check w$ 就是这个 $2$-丛的过渡数据。具体地说，$\check w$ 记录的是沿三重交叠把局部提升粘起来时出现的符号偏差：它不必是上闭链，其 coboundary 正是障碍类 $\alpha$，这也解释了为何在扭曲的图景里只能要求提升“相差 $\check w$”而不是严格相容。[来源：nLab "twisted spin structure"，https:/\allowbreak /\allowbreak ncatlab.org/\allowbreak nlab/\allowbreak show/\allowbreak twisted+\allowbreak spin+\allowbreak structure（2026-11-15 访问）]}

另一种描述较不一般，但借助先前讨论的材料也许更易理解。假设有光滑超曲面 $\Delta\subset B$，其上有实线丛 $\beta$，而 $w$ 由 $w_1(\beta)\in H^1(\Delta;\Z/2)$ 经推前得到。为简化叙述，再假设 $\Delta$ 有余定向。沿 $\Delta$ 切开 $B$，得到有两个边界部分 $\partial_+\widetilde B,\allowbreak \partial_-\widetilde B\cong\Delta$ 的流形 $\widetilde B$。此时，$F$ 上的扭曲 Pin 结构，就是拉回丛 $\widetilde F\to\widetilde B$ 上的 Pin 结构 $\widetilde P^\#$，连同同构 $\widetilde P^\#|\partial_+\widetilde B\cong\widetilde P^\#|\partial_-\widetilde B\otimes\beta$，它覆盖显然的识别 $\widetilde F|\partial_-\widetilde B=F|\Delta=\widetilde F|\partial_+\widetilde B$。
\end{note}

Pin 与 $\bar\partial$ 算子指标理论的关联，源自它和 $w_2$ 的关系。为看一个基本例子，暂时回到引理 \sr{st:11.7}{11.7}、\sr{st:11.12}{11.12} 的情形。考虑环路 $\rho\in\mathcal L_k\operatorname{Gr}(V)$、相应向量丛 $F_\rho\to S^1$，以及圆盘上的算子 $D_{D,\allowbreak \rho}$。令 $\widetilde F_\rho$ 为 $F_\rho$ 沿端点识别映射 $[0;1]\to S^1$ 的拉回。希望赋予它
\begin{equation}\label{eq:11.32}
\text{Pin 结构 }\widetilde P_\rho^\#,\quad\text{以及同构 }\widetilde P_1^\#\cong\widetilde P_0^\#\otimes\lambda^{\mathrm{top}}(\rho(1))^{\otimes k}.
\end{equation}
当 $k$ 为偶数时，这等价于 $F_\rho$ 上的一个 Pin 结构。当 $k$ 为奇数时，它是\stref{笔记}{11.16}中扭曲 Pin 结构定义的特例，其中 $F=F_\rho$，超曲面为点 $\Delta=\{0\}\subset S^1$，$\beta\to\{0\}$ 就是实直线 $\lambda^{\mathrm{top}}(\rho(1))$。这显然可推广到环路族，此时\eref{11.32}的存在障碍来自类
\begin{equation}\label{eq:11.33}
T(w_2)+kU(\mu)\in H^1(\mathcal L_k\operatorname{Gr}(V);\Z/2).
\end{equation}
\begin{lemma}\label{st:11.17}
$F_\rho$ 上一个形如\eref{11.32}的扭曲 Pin 结构的选择，确定同构 $\det(D_{D,\allowbreak \rho})\cong\lambda^{\mathrm{top}}(\rho(0))$。
\end{lemma}
\begin{proof}
对每个 $\rho$，扭曲 Pin 结构有两个可能的同构类；它们构成 $\mathcal L_k\operatorname{Gr}(V)$ 的一个二重覆盖，由\eref{11.33}分类。由\stref{引理}{11.7}，它同构于纤维为 $\det(D_{D,\allowbreak \rho})^\vee\otimes\lambda^{\mathrm{top}}(\rho(0))$ 的线丛所对应的覆盖。因此，要固定两个丛间的同构，只需考察它们在 $\mathcal L_k\operatorname{Gr}(V)$ 的某个非空连通子集上的限制。

这与\stref{引理}{11.12}的策略十分相似，事实上 $k=0$ 的情形只是该证明的改写。考虑常值环路 $\rho$，此时 $F_\rho$ 平凡，因而其上有选定的平凡 Pin 结构。规则是：若 Pin 结构平凡，就选择由 $\ker(D_{D,\allowbreak \rho})\cong\rho(0)$ 诱导的同构 $\det(D_{D,\allowbreak \rho})\cong
\SourcePage{177}
\lambda^{\mathrm{top}}(\rho(0))$；若 Pin 结构非平凡，就反转其符号。

再假设 $k=-1$。取正交复分解 $V=V_+\oplus V_-$，其中 $V_-$ 的复维数为一，考虑形如 $\rho(s)=\Lambda_+\oplus\exp(-\pi sI_V)\Lambda_-$ 的环路，其中 $\Lambda=\Lambda_+\oplus\Lambda_-$ 为相应分解的 Lagrangian 子空间。$F_\rho$ 上有一个显然的联络，其单值变换为 $A=\Id_{\Lambda_+}\oplus(-\Id_{\Lambda_-})$。一个扭曲 Pin 结构\eref{11.32}的选择，等价于将 $A$ 提升为 Pin 结构同构 $\widetilde P_0^\#\otimes\lambda^{\mathrm{top}}(\rho(0))\to\widetilde P_1^\#$。运用\eref{11.31}可知，它对应于 $A\otimes\Id$ 在 $\operatorname{Pin}(\Lambda\otimes\lambda^{\mathrm{top}}(\Lambda))\times_{\Z/2}S(\lambda^{\mathrm{top}}(\Lambda))$ 中的一个原像。由\stref{引理}{11.14}，两处 $\lambda^{\mathrm{top}}(\Lambda)$ 抵消，所以这归结为 $\Lambda_-$ 的一个定向。另一方面，将算子分解为 $V_\pm$ 两部分，并运用\stref{引理}{11.7}证明的 (iii)，可见 $D_{D,\allowbreak \rho}$ 为满射，其核同构地映到 $\Lambda_+$。因此，$\Lambda_-$ 的定向典范确定同构 $\det(D_{D,\allowbreak \rho})\cong\lambda^{\mathrm{top}}(\Lambda)$，正如所需。它对分解 $V=V_+\oplus V_-$ 和 $\Lambda=\Lambda_+\oplus\Lambda_-$ 的选择连续。

对一般的 $k$，考虑环路 $\rho(s)=\Lambda_+\oplus\exp(k\pi sI_V)\Lambda_-$。论证中 Pin 结构的部分只依赖 $k$ 的奇偶性，因此与之前相同。在行列式线丛方面，可用粘合公式\eref{11.11}得到典范同构 $\det(D_{D,\allowbreak \rho})\cong\lambda^{\mathrm{top}}(\Lambda_+)\otimes(\Lambda_-)^{\otimes(1+k)}$。
\end{proof}
